\documentclass[a4paper,10pt]{article}
\usepackage[francais]{babel}
\usepackage{url}

\selectlanguage{francais}
\input{letter-icube-uds.sty}

\selectlanguage{francais}
\begin{document}

\name{St\'ephane \textsc{Genaud}}
\fonction{Professeur des Universités}
\telephoneillkirch{+33 (0)3688 54542}
\email{genaud@unistra.fr}


\begin{letter}[Recommandation d'Etienne Michon pour la qualification aux
  fonctions de Maître de Conférences]{à}{CNU 27}

  Je connais  Etienne depuis  son intégration  dans notre  équipe à  travers son
  stage de  recherche Master, co-encadré  par Julien Gossa  et moi même,  sur le
  thème de l'allocation et l'ordonnancement  de ressources dans les clouds IaaS,
  Dans son  travail de thèse  à suivre, Etienne  a beaucoup contribué  au projet
  \emph{Schlouder}. Ce  logiciel est  un courtier  côté-client de  ressources de
  cloud IaaS.  Il accepte en  entrée des stratégies globales d'ordonnancement et
  d'allocation dynamiques multi-critères de  tâches indépendantes ou de workflow
  (souvent rapidité d'exécution versus le coût économique).  Cet outil permet de
  mener des  expériences dans  des environnements  réels de  clouds (s'interface
  avec OpenStack ou  EC2).  En dernière partie de sa  thèse et aujourd'hui alors
  qu'il occupe un  poste d'ATER à plein temps, Etienne  travaille à proposer une
  prédiction  des  résultats des  stratégies  d'ordonnancement  choisies par  le
  client.  Cette prédiction  est obtenue à partir d'une  simulation à évènements
  discrets basée sur le logiciel \textsc{SimGrid}.

  Etienne a soutenu sa thèse de doctorat  le 5 juin 2015.  Son travail durant sa
  thèse a fait avancer de manière considérable notre réflexion, en particulier à
  travers les expériences scientifique menées par Etienne. Pour mettre en oeuvre
  les idées générales  de l'équipe, Etienne a fourni un  travail considérable et
  rigoureux de  développement et  d'analyse qui ont  amené les  réalisations que
  nous  décrivons dans  un  article  de revue  actuellement  en soumission  dans
  \textit{Future Generation Computer Systems}.  Le processus de première lecture
  de cet  article à été  particulièrement long (13 mois)  et une révision  a été
  soumise récemment.

  D'un point de vue humain, Etienne a sans conteste de grandes qualités d'écoute
  et de  rigueur qui sont été  extrêment bénéfiques dans un  travail d'équipe de
  recherche. Il sait  prendre en compte les demandes et  les interrogations dans
  le cadre  de sa  recherche et les  traiter de manière  systématique. Il  a été
  durant sa thèse  une personne incontournable pour l'avancement  de nos travaux
  de recherche.

  Pour toutes ces raisons, je recommande  Etienne Michon pour la qualification
  aux fonctions de Maître de conférences en section CNU 27.

\end{letter}
\begin{flushright}
\includegraphics[width=.20\textwidth]{signgenaud.jpg}
\end{flushright}
\end{document}
